\subsection{A norm-one list with a matched finite upper bound}
\label{sec:norm-one-direct-list}

The next example concerns the size of one list, rather than the number
of exceptional parameters on a received line. Its incidence parameters
are those of the classical projective-plane design, with each coordinate
doubled. We make no novelty claim for that geometry.

\begin{theorem}\label{thm:norm-one-direct-list}
Let $p\ge53$ be prime, let $E=\mathbb F_{p^3}$, and put
\[
 L=p^2+p+1,\qquad N=2L,\qquad
 G=\ker\operatorname{Norm}_{E/\mathbb F_p},\qquad
 D=\{x\in E^*:x^2\in G\}.
\]
The dimension-three Reed--Solomon code on $D$ has a received word
$f(x)=x^{2p+2}$ whose maximum agreement is $A=2p+2$ and whose complete
list at that threshold consists of the $L=N/2$ polynomials
\[
 Q_a(X)=aX^2-a^{p^2+1},\qquad
 \operatorname{Norm}_{E/\mathbb F_p}(a)=-1.
\]
This list attains the Johnson counting bound exactly. At the lower
threshold
\[
 T=\lceil19p/10\rceil
\]
the complete list of $f$ is still exactly the displayed $N/2$ witnesses,
and the same code satisfies
\[
 N a_1(3/N)<T<\sqrt{2N},
\]
\[
 \frac N2\le\max_{w\in E^D}
 \#\{Q\in E[X]:\deg Q\le2,\ |\{x\in D:w(x)=Q(x)\}|\ge T\}
 \le22N.
\]
Thus the lower bound is within a factor $44$ of an explicit uniform
finite first-order list bound at the same parameters. Moreover
$T/\sqrt N\to19/(10\sqrt2)$, strictly between the limiting
first-order value $\sqrt{3/2}$ and the Johnson value $\sqrt2$.
Both margins therefore have order $\sqrt N$. Here the rate tends
to zero, $p=\Theta(\sqrt N)$, and $|E|=\Theta(N^{3/2})$; this is not a
fixed-rate or prime-alphabet assertion.
\end{theorem}

\begin{proof}
The group $G$ has odd order $L$ and lies in the squares of $E^*$, so
$|D|=2L$. For $\operatorname{Norm}(a)=-1$, substitute
$y=a^{p^2}z$ in $y^{p+1}=ay-a^{p^2+1}$. The result is
\[
 z^{p+1}-z+1=0,\qquad z^p=1-1/z.
\]
Neither $0$ nor $1$ is a root. The fractional linear map $1-1/z$
has third iterate the identity, so all roots lie in $E$. The polynomial
is squarefree: its derivative is $z^p-1$, and a common zero would leave
residual $1$. Its $p+1$ roots satisfy
$z^{p^2}=1/(1-z)$ and hence $\operatorname{Norm}(z)=-1$.
Consequently all $p+1$ values of $y$ have norm one. Each has two square
roots in $D$, proving exactly $A$ matches for each of the $L$ distinct
quadratics displayed above.

For any quadratic $Q$, the polynomial $f-Q$ has degree $A$, so no
quadratic has more than $A$ matches. If the complete threshold-$A$ list
has size $m$, let $r_x$ count its matching witnesses at $x$. Distinct
quadratics share at most two matching coordinates. Therefore
\[
 \sum_x r_x=mA,\qquad
 \sum_xr_x^2\le mA+2m(m-1).
\]
Cauchy's inequality gives
\[
 m\le\frac{N(A-2)}{A^2-2N}=L,
 \qquad A^2-2N=4p.
\]
This proves completeness and exact Johnson saturation. Equality also
shows that every coordinate belongs to $p+1$ supports and every pair
of supports shares two coordinates.

We next show that every quadratic outside this bank has at most $p+1$
matches. We use the following elementary semilinear observation. For
$v\ne0$, a nonzero root of $z^{p+1}=uz+v$ satisfies
$z^p=u+v/z$. The threefold Frobenius-twisted composition of this
fractional linear map is the identity exactly when
$u\ne0$ and $v=-u^{1+p^2}$: this follows by multiplying the matrices
$\left(\begin{smallmatrix}u^{p^i}&v^{p^i}\\1&0\end{smallmatrix}\right)$
for $i=2,1,0$. Otherwise it has at most two fixed points.
The identity case has exactly $p+1$ roots by the same squarefreeness
and Frobenius argument as above. When $v=0$, there is at most one
nonzero root.

For an even quadratic $Q=aX^2+c$, apply this observation to
$y=x^2$. More than four matches requires
$a\ne0$ and $c=-a^{1+p^2}$. All its $y$-roots have norm
$-\operatorname{Norm}(a)$, so any match on $D$ then requires
$\operatorname{Norm}(a)=-1$.
For $Q=aX^2+bX+c$ with $b\ne0$, put
\[
 r(X)=1-a^pQ(X)-c^pX^2.
\]
At a matching $x\in D$, the norm-one condition gives
$x^2Q(x)^p=1$, hence $r(x)=b^px^{p+2}$. Thus
\[
 r(X)^2-b^{2p}X^2Q(X)
\]
vanishes at every match. If this polynomial is nonzero, there are at
most four matches. If it is zero, $X$ divides $r$ and
\[
 H(X)=\frac{r(X)}{b^pX}
\]
is a polynomial of degree at most one, with $H^2=Q$.
The original unsquared identity forces
$x^{p+1}=H(x)$ at every match, giving at most $p+1$ matches.
As $p+1\ge4$, every outsider has at most $p+1$ matches.
Consequently the displayed bank is the complete list at every integer
threshold $p+1<T'\le2p+2$, including $T$.

For the uniform upper bound at $T$, use the derivative-weighted support
of~\cite[Proposition~5.10]{dkt2026} with
\[
 m=6,\quad B_\partial=3,\quad B=B_{\rm jet}=3T,
 \quad H=B_Z=40B.
\]
Its coefficient count is $G_{\rm coef}=4B^2-2B$ and its six local-rank
contributions are $4,8,12,15,14,9$, totaling $R=62$.
For $p\ge53$, $57p/10\le B\le6p$, and
\[
 39G_{\rm coef}-40NR
 \ge\frac{2711}{25}p^2-5428p-4960>0.
\]
Thus $G_{\rm coef}>NR$ and the graded support surplus is at least
$B(39G_{\rm coef}-40NR)+(G_{\rm coef}-NR)>0$.
This certificate applies uniformly to evaluation domains and received
words. Reconstruction requires only $p>\max(2,3)$.

In~\cite[Eq.~(54), Lemma~5.11]{dkt2026}, message degree $2$ and
$(B,M)=(3T,3)$ give $F_{\rm reg}=7B-12$ and $S=5B-9$.
Writing $\lambda=(N-2)/(T-2)$, the fixed-word bound is therefore
\[
 \lambda F_{\rm reg}+S
 =21N+30\lambda+15T-51\le21N+90p-51\le22N.
\]
Here $T-2\ge p+1$ gives $\lambda\le2p$, and
$N-90p+51=2p^2-88p+53>0$.
The exhibited complete list supplies the lower bound at $T$.
Finally, $T<2p<\sqrt{2N}$, while the first-order curve estimate gives
\[
 N a_1(3/N)\le\sqrt{3N/2}+(3N/8)^{1/4}
 <\frac74p+\sqrt p<T
\]
using $T\ge19p/10$, $\sqrt p>20/3$, and $p\ge53$.
\end{proof}

The matched comparison is codewide, not just a bound along one received
line. Its normalized threshold and normalized separation from the
first-order curve both vanish with the rate; it makes no claim of
constant-factor tightness for exceptional-label or common-agreement
bounds.


\begin{corollary}\label{cor:norm-one-full-interior-band}
With the code and word of Theorem~\ref{thm:norm-one-direct-list}, let
$p\ge401$. At every integer threshold
\[
 \lceil\sqrt3\,p\rceil\le T\le2p+2,
\]
the displayed word has exactly $N/2$ nearby quadratics, and the maximum
list size over all received words is at most $137N$.
Consequently, for every fixed $c\in(\sqrt3,2)$, at threshold
$T=\lceil cp\rceil$ the maximum list size is $\Theta(N)$.
For sufficiently large $p$ these thresholds lie strictly between the
first-order curve and the Johnson threshold, with both margins of
order $\sqrt N$.
\end{corollary}

\begin{proof}
The classification in the theorem already gives the exact $N/2$ list
throughout this band. For the uniform upper bound, take
\[
 m=16,\qquad B_\partial=8,\qquad B=B_{\rm jet}=8T,
 \qquad H=B_Z=100B.
\]
The derivative-weighted coefficient count and local rank are
\[
 G=9B^2-27B,\qquad R=852.
\]
Since $8\sqrt3\,p\le B\le16p+16$, we have
\[
 99G-100NR
 \ge672p^2-213168p-213168>0
 \qquad(p\ge401).
\]
The graded surplus $(H-B+1)G-(H+1)NR$ is therefore positive.
The characteristic guard is $p>8$.
Equation~(54) and Lemma~5.11 of~\cite{dkt2026} give
$F_{\rm reg}=17B-72$ and $S=15B-64$. With
$\lambda=(N-2)/(T-2)\le2p$, the list bound is
\[
 \lambda F_{\rm reg}+S
 =136N+200\lambda+120T-336
 \le136N+640p-96\le137N.
\]
Finally, $Na_1(3/N)=\sqrt3\,p+O(\sqrt p)$ and
$\sqrt{2N}=2p+O(1)$. Thus every fixed $c\in(\sqrt3,2)$
has the asserted two strict margins for all sufficiently large $p$.
The finite band in the first statement need not itself lie wholly
between those two boundaries.
\end{proof}
